\documentclass[conference]{IEEEtran}
\IEEEoverridecommandlockouts
\usepackage{cite}
\usepackage{amsmath,amssymb,amsfonts}
\usepackage{graphicx}
\usepackage{textcomp}
\usepackage{xcolor}
\usepackage{booktabs}
\usepackage{array}
\usepackage{tabularx}
\usepackage[hidelinks]{hyperref}
\graphicspath{{figures/}{../../../paper/generated/}}
\setlength{\textfloatsep}{7pt plus 2pt minus 2pt}
\setlength{\floatsep}{7pt plus 2pt minus 2pt}
\setlength{\intextsep}{7pt plus 2pt minus 2pt}
\renewcommand{\arraystretch}{0.98}
\newcommand{\AuditTotalRows}{1257}
\newcommand{\AuditValidationRows}{140}
\newcommand{\AuditValidationExcluded}{4}
\newcommand{\AuditTestCandidateRows}{343}
\newcommand{\AuditTestExcluded}{10}
\newcommand{\AuditTestRows}{333}
\newcommand{\AuditTailRows}{73}
\newcommand{\AuditSelectedModel}{ExtraTrees}
\newcommand{\AuditSelectedFamily}{weather only}
\newcommand{\AuditSelectedRtwo}{-0.014}
\newcommand{\AuditSelectedRMSE}{0.669}
\newcommand{\AuditSelectedMAE}{0.474}
\newcommand{\AuditSelectedDeltaRMSE}{-0.005}
\newcommand{\AuditSelectedDeltaHigh}{0.009}
\newcommand{\AuditGateACI}{[-0.019, 0.009]}
\newcommand{\AuditGateBOneDeltaRMSE}{-0.012}
\newcommand{\AuditGateBOneCI}{[-0.029, 0.002]}
\newcommand{\AuditGateBTwoCI}{[-0.019, 0.009]}
\newcommand{\AuditTailGate}{FAIL}
\newcommand{\AuditPositiveStates}{4/12}
\newcommand{\AuditLeakageJaccard}{0.605}
\newcommand{\AuditLeakageRank}{0.942}


\title{Auditing Weather-Feature Reliance in Detrended Crop-Yield Models}
\author{\IEEEauthorblockN{Anonymous Submission}}

\begin{document}
\maketitle

\begin{abstract}
Interpretable machine learning cannot support substantive claims about below-trend
crop-yield events unless its fitted predictive model has adequate out-of-sample
fidelity. We audit this condition in a U.S. crop--state--year panel with
\AuditTotalRows{} rows for Barley, Canola, Oats, and Wheat from 1990--2025.
Trends and residual scales are fit on training years only; selection uses
2012--2015 validation, and 2016--2025 remains locked until final evaluation. The
validation-selected \AuditSelectedModel{} \AuditSelectedFamily{} model has final
$R^2=\AuditSelectedRtwo{}$, RMSE \AuditSelectedRMSE{} t ha$^{-1}$, and paired
$\Delta$RMSE versus a zero-residual baseline of \AuditSelectedDeltaRMSE{} with
upper 95\% bound \AuditSelectedDeltaHigh{}. The selected model fails Gate A, and
weather features also fail the primary incremental-value Gate B1. We therefore
retain SHAP, LIME, permutation, ablation, and response-curve results only as
descriptions of a fitted function, not evidence that weather caused observed
events. The contribution is a reproducible stop-gate protocol with paired
year-block uncertainty, tail checks, robustness checks, and an explicit warning
that full-series detrending is retrospective only.
\end{abstract}

\begin{IEEEkeywords}
crop yield, explainable machine learning, temporal validation, negative results,
reproducibility
\end{IEEEkeywords}

\section{Introduction}
Machine-learning yield forecasts are increasingly built from weather, remote
sensing, management, and crop-model variables. Studies report useful predictive
systems across large and small panels, diverse model families, and hybrid
formulations \cite{Paudel2021,Meroni2021,Khaki2019,LengHall2020,Klompenburg2020,Chlingaryan2018,Shahhosseini2021}.
Yet a plausible feature attribution is not causal evidence, and agreement between
explanation methods cannot compensate for a model that does not predict the future
target reliably. This distinction is especially important for below-trend yield
events, where a compelling explanation of an individual fitted prediction can be
mistaken for an explanation of the observed loss.

Weather extremes and climate variability are materially associated with crop-yield
variation at several spatial scales \cite{Lesk2016,Ray2015,Lobell2011,SchlenkerRoberts2009,Zampieri2017,Vogel2019,Heino2023,Schierhorn2021,Sjulgard2023}.
Those findings motivate careful predictive evaluation; they do not license an
event-level attribution from this restricted panel. We ask a narrower question:
does a leakage-safe residual model have enough future-period fidelity and stability
to permit substantive weather-feature-reliance interpretation? The answer is
evaluated before XAI is interpreted.

The paper makes three contributions. First, it uses train-only target construction
and validation-then-locked-test selection. Second, it specifies a quantitative
fidelity gate with same-task baselines, paired year-block intervals, tail metrics,
and pre-specified robustness checks. Third, it contrasts a prospective target with
a retrospective full-series detrending diagnostic. A gate failure is a negative
methodological result: it stops causal or event-driver claims, rather than choosing
the most attractive explanation.

\section{Related Work}
\subsection{Yield Modeling and Detrending}
Detrending separates a local historical yield trajectory from residual variation,
but the result can depend on the detrending procedure and information available at
the evaluation date \cite{Lu2017,Meng2024}. We therefore fit each crop--state trend
and residual scale on training observations only. Ridge regression, Random Forest,
and ExtraTrees provide deliberately distinct linear and tree-based model families
\cite{HoerlKennard1970,Breiman2001,Geurts2006}; the implementations use
scikit-learn \cite{Pedregosa2011}. This is a finite pre-specified comparison, not
a search over the locked test period.

\subsection{Explanation Methods and Their Boundary}
SHAP and LIME explain a model output under their respective approximation and
reference choices \cite{LundbergLee2017,Ribeiro2016}. Permutation-style analyses,
model reliance, accumulated local effects, and conditional importance likewise
describe feature dependence or behavior of a fitted predictor
\cite{Fisher2019,ApleyZhu2020,Strobl2008}. They can be useful diagnostic artifacts.
However, none turns failed prospective predictive fidelity into evidence about an
observed yield anomaly. The gate below formalizes this boundary.

\subsection{Evaluation and Reproducibility}
Random splits can be inappropriate when temporal or hierarchical structure is
present \cite{Roberts2017}. Leakage can create seemingly strong scientific
machine-learning results while invalidating the intended forecast claim
\cite{KapoorNarayanan2023}. Reproducible ML also requires explicit code, data
paths, settings, and numerical evidence \cite{Pineau2021}. Accordingly, every
manuscript number is generated from an auditable run: configuration, prediction
vectors, bootstrap draws, citation check, and final PDF are all versioned outputs.

\section{Data and Protocol}
\subsection{Panel, Features, and Target}
The unit is a crop--state--year row. The processed panel contains \AuditTotalRows{}
observations from 1990--2025 for Barley, Canola, Oats, and Wheat across 12 U.S.
states. Yield in t ha$^{-1}$ comes from USDA NASS Quick Stats
\cite{USDANASSQuickStats}; daily maximum/minimum temperature, precipitation, and
shortwave radiation come from NASA POWER \cite{NASAPOWER2026}. The primary design
has 35 full-season weather features. A separate sensitivity design adds early-,
middle-, and late-season proxies without redefining the target or locked population.

\begin{table}[t]
\centering
\caption{Data-flow and final-test eligibility reconstructed from the raw inputs.}
\label{tab:dataflow}
\scriptsize
\resizebox{\columnwidth}{!}{\begin{tabular}{lrl}
\toprule
Stage & Rows & Rule \\
\midrule
raw yield file & 1291 & source table supplied in repository \\
analysis year filter & 1257 & exclude 34 raw 1989 rows because analysis starts in 1990 \\
processed model frame & 1257 & exact reconstruction from raw yields plus NASA POWER daily input \\
final test eligible & 333 & 2016-2025 rows with at least three prior series observations \\
\bottomrule
\end{tabular}
}
\end{table}

For a test interval, each crop--state linear trend and residual scale is estimated
using only its associated training rows. A row is eligible only with at least three
training observations. The validation period contains \AuditValidationRows{} eligible
rows after \AuditValidationExcluded{} pre-specified exclusions. The locked test
starts with \AuditTestCandidateRows{} candidates, excludes \AuditTestExcluded{},
and evaluates \AuditTestRows{} rows. Future yields cannot change a previously
constructed evaluation target; this property has a unit test and row-level audit.
The learned target is the raw train-only residual in t ha$^{-1}$; the standardized
residual $z=r/\hat{\sigma}_{\mathrm{train}}$ defines below-trend events only.
This distinction prevents a pooled raw-error score from being misread as a
standardized event-recovery metric. Min-history rules of 3, 5, 8, and 10 years are
reported as sensitivity analyses because both eligibility and the fitted residual
scale can change with available history. Three observations are the minimum that
leaves one positive residual degree of freedom for estimating a training residual
scale; this is not claimed to provide a stable series-specific variance estimate.

\subsection{Selection and Baselines}
Ridge $\alpha\in\{1,10\}$ and Random Forest/ExtraTrees with 160 trees and
$\mathrm{min\_samples\_leaf}\in\{1,2\}$ are evaluated on metadata-only,
weather-only, and full representations. Stochastic fits use fixed seeds
$\{7,17,27,37,47\}$; numeric imputation/scaling, categorical encoding, and model
fitting occur within the training fold. Mean per-row seed-aggregated validation
RMSE chooses the configuration, with lexicographic tie-breaking and no locked-test
access. Seed-level RMSE standard deviations are reported so tiny differences are
not presented as deterministic wins.

The final audit compares learned predictions on identical rows to zero residual,
training-mean residual, and crop training-mean residual baselines. Prediction
vectors are hashed and pairwise maximum absolute differences are retained, so any
nearly coincident baselines are demonstrated rather than assumed. We use 2,000
deterministic calendar-year block bootstrap replicates and percentile 95\%
intervals. An invalid replicate fails the run; none are silently dropped.
We do not implement the stationary bootstrap itself; we cite it for the
dependent-resampling motivation and resample complete calendar years as paired
blocks \cite{PolitisRomano1994}. State-year and crop--state cluster schemes are
reported as pre-specified sensitivity analyses.

\begin{figure*}[t]
  \centering
  \includegraphics[width=0.96\textwidth]{figure_audit_workflow.png}
  \caption{Prospective audit workflow. Selection is complete before the locked
  2016--2025 test; the final fidelity gate controls whether interpretation may
  proceed.}
  \label{fig:workflow}
\end{figure*}

\subsection{Pre-specified Fidelity Gate}
For a paired comparison, model-minus-baseline error passes only when the upper
95\% year-block interval is below zero. Gate A is predictive fidelity of the
validation-locked choice. Gate B1 is the primary incremental-weather-value test:
the validation-selected Weather only versus Metadata only contrast. Gate B2
(Full versus Metadata only) is sensitivity only and never triggers re-selection.
Weather-specific claims require both Gate A and Gate B1.
The primary operational tail is $z<-1$;
$z<-1.5$ and $z<-2$ are sensitivity analyses because their smaller samples reduce
power. Sign accuracy on a negative-only subset is reported only as a diagnostic,
not a gate component. Instead, rank recovery requires a positive lower bootstrap
bound and a within-year permutation $p<.05$. Top-$k$ recovery reports random
expectation $k/n$, lift, hypergeometric probability, permutation $p$, and a
year-block lift interval. Thus a severe-event row cannot pass simply by overlapping
one arbitrary event. The fixed policy is in Table~\ref{tab:gate}.

\begin{table*}[t]
\centering
\caption{Fixed gate policy. Gate A and primary Gate B1 are independent; Gate B2
is sensitivity only. Machine-readable values are in
\nolinkurl{configs/fidelity_gate.yaml}.}
\label{tab:gate}
\footnotesize
\renewcommand{\arraystretch}{1.12}
\begin{tabular}{p{0.16\textwidth}p{0.24\textwidth}p{0.34\textwidth}p{0.16\textwidth}}
\toprule
Module & Question & Pass condition & Claim scope \\
\midrule
A -- Overall adequacy & Does the locked model improve on the same-target baseline? & Upper paired 95\% CI of $\Delta$RMSE $<0$. & Overall predictive claim \\
B -- Weather value & Does Full improve on Metadata-only on the same locked rows? & Upper paired 95\% CI of $\Delta$RMSE $<0$. & Weather-specific predictive reliance \\
E -- Event recovery & Does the model recover the primary below-trend tail beyond null checks? & Tail RMSE/MAE, rank, permutation, and top-$k$ checks all pass. & Observed below-trend event claim \\
D -- Diagnostic & How does Weather-only compare with Metadata-only? & Reported only; not a decision module. & Descriptive diagnostic only \\
\bottomrule
\end{tabular}

\end{table*}

\begin{center}
\fbox{\begin{minipage}{0.94\linewidth}\scriptsize
\textbf{Algorithm 1: Locked fidelity audit}\par
\textbf{Input:} panel, fixed split/configuration, seed list, gate configuration.\par
1. Fit detrending and preprocessing on each training period only.\par
2. Select one configuration on validation RMSE; lock it.\par
3. Refit the locked choice through 2015 and predict 2016--2025.\par
4. Compare same-row baselines with paired year-block bootstrap intervals.\par
5. Evaluate null-aware primary-tail rank/top-$k$ recovery and sensitivity tails.\par
6. Return Gate A, primary Gate B1, and sensitivity Gate B2 separately.
Substantive observed-event claims require Gate A, while weather-specific claims
require both Gate A and Gate B1. Gate B2 cannot trigger model re-selection or
replace Gate B1.
\end{minipage}}
\end{center}

\section{Results}
\subsection{Validation and Locked Test}
Table~\ref{tab:selection} shows leading validation configurations. The selected
model is \AuditSelectedModel{} with \AuditSelectedFamily{} features. The complete
grid, seed-level metrics, tie rule, and locked-test access flag are supplied in the
artifact. A non-selected ExtraTrees full model can have a small final point estimate
without changing the decision: it was not validation-selected, and the locked test
is never used to switch configurations.

\begin{table}[t]
\centering
\caption{Validation selection on 2012--2015. RMSE is computed after fixed-seed
prediction aggregation; Seed SD is the across-seed RMSE standard deviation.}
\label{tab:selection}
\scriptsize
\resizebox{\columnwidth}{!}{\begin{tabular}{lllrrrccc}
\toprule
Model & Features & Role & $n$ & RMSE & Seed SD & $P$(rank 1) & 1-SE & Selected \\
\midrule
ExtraTrees (leaf=1) & Weather only & Module A; D & 140 & 0.384 & 0.003 & 0.60 & Yes & Yes \\
ExtraTrees (leaf=2) & Weather only & Validation candidate & 140 & 0.385 & 0.001 & 0.20 & No & No \\
Random Forest (leaf=1) & Weather only & Validation candidate & 140 & 0.387 & 0.003 & 0.20 & No & No \\
Random Forest (leaf=2) & Weather only & Validation candidate & 140 & 0.387 & 0.002 & 0.00 & No & No \\
ExtraTrees (leaf=2) & Full & Validation candidate & 140 & 0.389 & 0.003 & 0.00 & No & No \\
Random Forest (leaf=1) & Full & Validation candidate & 140 & 0.390 & 0.002 & 0.00 & No & No \\
ExtraTrees (leaf=1) & Full & Module B Full & 140 & 0.391 & 0.003 & 0.00 & No & No \\
Ridge ($\alpha=10$) & Metadata only & Best Ridge candidate & 140 & 0.407 & N/A & N/A & No & No \\
ExtraTrees (leaf=1) & Metadata only & Module B/D Metadata & 140 & 0.407 & 0.000 & 0.00 & No & No \\
\bottomrule
\end{tabular}
}
\end{table}

The selected weather-only ExtraTrees configuration ranks first in 3/5 stochastic
validation seeds. Its advantage is small relative to several seed standard
deviations, so the selected row is a reproducible protocol choice rather than a
claim of decisive superiority. The one-standard-error flag and every seed-level
score are retained; no final-test statistic changes this selection.

\begin{table}[t]
\centering
\caption{Validation stability of the leading configurations. The rank probability
is across fixed stochastic seeds; 1-SE denotes inclusion under the pre-specified
one-standard-error diagnostic, not a post-hoc replacement.}
\label{tab:stability}
\scriptsize
\resizebox{\columnwidth}{!}{\begin{tabular}{llrrrr}
\toprule
Model & Features & Mean RMSE & SD & P(rank=1) & 1-SE \\
\midrule
ExtraTrees (leaf=1) & Weather only & 0.384 & 0.003 & 0.60 & Yes \\
ExtraTrees (leaf=2) & Weather only & 0.385 & 0.001 & 0.20 & No \\
Random Forest (leaf=1) & Weather only & 0.387 & 0.003 & 0.20 & No \\
Random Forest (leaf=2) & Weather only & 0.387 & 0.002 & 0.00 & No \\
ExtraTrees (leaf=2) & Full & 0.389 & 0.003 & 0.00 & No \\
Random Forest (leaf=1) & Full & 0.390 & 0.002 & 0.00 & No \\
\bottomrule
\end{tabular}
}
\end{table}

On the locked test, the selected model has $R^2=\AuditSelectedRtwo{}$ and RMSE
\AuditSelectedRMSE{} t ha$^{-1}$. Its paired RMSE difference from zero crosses
zero, so Gate A fails. Table~\ref{tab:baselines} includes all same-task naive
baselines and the selected configuration across feature families. Figure
\ref{fig:finalaudit} separates three locked-test questions. Gate A evaluates
whether the validation-selected weather-only model improves on the zero-residual
baseline. Gate B1, the primary incremental-weather test, compares the same
selected weather-only model with metadata-only context. Gate B2 compares the full
and metadata-only representations as a sensitivity analysis; it cannot trigger
re-selection or replace Gate B1. All comparisons use paired predictions on
identical locked rows. Gate A, Gate B1, and Gate B2 all fail.

\begin{table*}[t]
\centering
\caption{Locked 2016--2025 same-task audit. Learned-model predictions are seed
aggregated. $\Delta$RMSE is model minus zero-residual RMSE with paired 95\%
year-block intervals; RMSE is t ha$^{-1}$.}
\label{tab:baselines}
\scriptsize
\resizebox{\textwidth}{!}{\input{../../../paper/generated/table_final_baselines.tex}}
\end{table*}

\subsection{Event-Detection Diagnostics}
Event detection is evaluated over all 333 locked rows, not only a tail preselected
by the observed outcome. The score is negative predicted residual and the score
threshold is fixed at zero before final evaluation. Table~\ref{tab:event} reports
PR-AUC, ROC-AUC, balanced accuracy, precision, recall, and F1. These are
diagnostics rather than Gate A components: they describe discrimination, but they
do not override the paired-error and null-aware rank failures. In particular, an
always-negative predictor cannot gain a gate pass through the negative-only sign
rate used in older drafts.

\begin{table*}[t]
\centering
\caption{Locked-test event-detection diagnostics on all rows. These values are
not tuned on the final test and do not replace Gate A components.}
\label{tab:event}
\scriptsize
\resizebox{\textwidth}{!}{\begin{tabular}{lrrrrrrr}
\toprule
Tail & Prev. & PR-AUC & ROC-AUC & Bal. acc. & Precision & Recall & F1 \\
\midrule
z$<$-1 & 0.219 & 0.338 & 0.641 & 0.589 & 0.294 & 0.548 & 0.383 \\
z$<$-1.5 & 0.156 & 0.294 & 0.672 & 0.611 & 0.228 & 0.596 & 0.330 \\
z$<$-2 & 0.108 & 0.217 & 0.672 & 0.598 & 0.154 & 0.583 & 0.244 \\
\bottomrule
\end{tabular}
}
\end{table*}

\begin{figure}[t]
  \centering
  \includegraphics[width=\linewidth]{figure_final_audit.png}
  \caption{Three distinct locked-test comparisons. Gate A evaluates predictive
  fidelity of the validation-selected weather-only model against the zero-residual
  baseline. Gate B1 is the primary incremental-weather comparison of the selected
  weather-only model against metadata-only context. Gate B2 compares the full and
  metadata-only models as a sensitivity analysis and cannot alter model selection.
  Values are paired $\Delta$RMSE estimates with 95\% calendar-year-block
  intervals; negative values favor the model on the left.}
  \label{fig:finalaudit}
\end{figure}

\subsection{Below-Trend Gate and Robustness}
The broad below-trend subset contains \AuditTailRows{} locked-test rows. Table
\ref{tab:tail} separates error, rank-null, and chance-adjusted event-recovery
components. The primary $z<-1$ top-10 lift is below random expectation and its
permutation result is not significant; rank uncertainty also crosses zero. Gate A
therefore fails without relying on a directional-sign statistic. Neither local
explanations nor global/group diagnostics are evidence of weather causes or event
drivers in this primary analysis.

\begin{table*}[t]
\centering
\caption{Tail components on train-only targets. Panel A reports error and rank
recovery; Panel B reports chance-adjusted Top-10 recovery. $H$ is the
hypergeometric probability and $P$ is the within-year permutation probability.
The Top-10 component passes only when lift $>1$, the lower lift CI $>1$, and both
$H$ and $P$ are $<.05$. Only $z<-1$ is primary; severe tails are sensitivities.}
\label{tab:tail}
\footnotesize
\renewcommand{\arraystretch}{1.08}
\textbf{Panel A: Error and rank recovery}\par\smallskip
\input{../../../paper/generated/table_tail_error_rank.tex}
\par\medskip
\textbf{Panel B: Chance-adjusted Top-10 recovery}\par\smallskip
\begin{tabular}{llrrrrrrl}
\toprule
Role & Tail & $n$ & Observed & Expected & Lift [95\% CI] & $H$ $p$ & $P$ $p$ & Status \\
\midrule
Primary & z$<$-1 & 73 & 1/10 & 1.37 & 0.73 [0.00, 2.82] & 0.794 & 0.793 & FAIL \\
Sensitivity & z$<$-1.5 & 52 & 1/10 & 1.92 & 0.52 [0.37, 2.19] & 0.907 & 0.935 & FAIL \\
Sensitivity & z$<$-2 & 36 & 1/10 & 2.78 & 0.36 [0.32, 1.68] & 0.979 & 0.985 & FAIL \\
\bottomrule
\end{tabular}

\end{table*}

\begin{table*}[t]
\centering
\caption{Gate decision components. Gate A controls substantive observed-event
claims; primary Gate B1 is additionally required for weather-specific claims;
Gate B2 is sensitivity only.}
\label{tab:gateab}
\scriptsize
\resizebox{\textwidth}{!}{\begin{tabular}{p{0.10\textwidth}p{0.22\textwidth}p{0.28\textwidth}p{0.24\textwidth}c}
\toprule
Gate & Role & Contrast & $\Delta$RMSE [95\% CI] & Status \\
\midrule
Gate A & primary overall fidelity & Selected Weather-only vs. zero residual & -0.005 [-0.019, 0.009] & FAIL \\
Gate B1 & primary conditional weather value & Full vs. Metadata-only & -0.012 [-0.029, 0.002] & FAIL \\
Gate B2 & diagnostic representation contrast & Weather-only vs. Metadata-only & -0.005 [-0.019, 0.009] & FAIL \\
\bottomrule
\end{tabular}
}
\end{table*}

Table~\ref{tab:robustness} gives exact temporal pass counts and summary ranges for
the pre-specified stage-feature, crop, spatial, and target-protocol checks. These
checks do not reverse the central decision; they show its stability boundary rather
than optimize around it.

\begin{table*}[t]
\centering
\caption{Robustness and protocol sensitivity. Exact criteria, row-level outputs,
and settings are retained in the audit artifact.}
\label{tab:robustness}
\footnotesize
\renewcommand{\arraystretch}{1.08}
\begin{tabular}{p{0.20\textwidth}p{0.32\textwidth}p{0.16\textwidth}p{0.24\textwidth}}
\toprule
Audit & Main result & Role & Interpretation \\
\midrule
Rolling origin & 3/4 negative; 1/4 CI pass & Diagnostic & Unstable \\
Stage proxies & $\Delta$RMSE -0.001 & Sensitivity & No material change \\
Crop-specific & $R^2$ -0.273 to 0.072 & Diagnostic & Unstable \\
Leave-one-state-out & 4/12 states improve & Diagnostic & Unstable \\
Full-series detrending & Jaccard 0.605 & Retrospective & Future yields enter target \\
\bottomrule
\end{tabular}

\end{table*}

\subsection{Target, Cluster, and Model-Family Sensitivity}
The three-observation rule is the pre-specified technical minimum, not an
empirically optimal or stable variance threshold. Eligibility is audited at each validation/final
cutoff. Histories 8 and 10 retain 313 rather than 333 final rows, so their favorable
intervals are sensitivity results from a different population and cannot alter
primary Gate A. Their eligible row IDs and target vectors are identical; their
predictions differ only at floating-point precision (maximum absolute difference
$1.11\times10^{-16}$), so the displayed metrics coincide. Table~\ref{tab:sensitivity}
separates target/model checks from resampling, so row counts and cluster counts
cannot be confused.

\begin{table*}[t]
\centering
\caption{Target and model sensitivity. History 3 is the primary technical minimum.
RMSE / zero gives model and zero-baseline RMSE. *Sensitivity PASS is evaluated on a
different eligible population and cannot alter the primary Gate A decision. History 8
and History 10 have identical eligible rows/targets; their maximum prediction
difference is $1.11\times10^{-16}$, so displayed metrics coincide.}
\label{tab:sensitivity}
\footnotesize
\renewcommand{\arraystretch}{1.05}
\input{../../../paper/generated/table_target_model_sensitivity.tex}
\end{table*}

\begin{table*}[t]
\centering
\caption{Paired resampling sensitivity on the locked 333-row population. The
number of clusters is a resampling property, never a row count.}
\label{tab:resampling}
\scriptsize
\resizebox{\textwidth}{!}{\begin{tabular}{lrrrlc}
\toprule
Resampling unit & Number clusters & Test rows & Delta RMSE & 95\% CI & Status \\
\midrule
year block & 10 & 333 & -0.005 & [-0.019, 0.009] & SENSITIVITY FAIL \\
state year block & 120 & 333 & -0.005 & [-0.022, 0.010] & SENSITIVITY FAIL \\
crop state cluster & 36 & 333 & -0.005 & [-0.033, 0.019] & SENSITIVITY FAIL \\
\bottomrule
\end{tabular}
}
\end{table*}

The standardized-target model is actually refit on the train-only residual $z$,
rather than obtained by rescaling raw-target predictions after the fact; its paired
interval also crosses zero. Test-period $z$ values need not have unit variance:
each crop--state series uses its training-only residual scale, and a small training
scale can yield larger future standardized residuals without leakage. Replacing ordinary least-squares detrending by a
training-only Huber trend changes the primary event count from 73 to 40, but its
paired interval still crosses zero. These analyses are deliberately reported as
target-construction sensitivities, not candidate routes to a favorable result.

\subsection{Temporal Stability and Capacity}
The four rolling-origin folds use non-overlapping future blocks after 2003, 2006,
2009, and 2012 training cutoffs. Three of four folds have negative point estimates;
only one of four has a 95\% CI fully below zero. A prefix learning curve retrains the same validation-locked
configuration using progressively longer historical prefixes, while retaining only
crop--state series with at least three prefix observations. None of the four
prefix intervals establishes improvement over zero. This makes the conclusion
conditional on the available panel and model family, not a claim that more data
cannot help.

\begin{table*}[t]
\centering
\caption{Exact temporal robustness and prefix learning-curve checks. Each entry
compares the locked weather-only configuration to zero residual with paired
calendar-year-block 95\% intervals.}
\label{tab:temporalcapacity}
\scriptsize
\resizebox{\textwidth}{!}{\begin{tabular}{lrrrlc}
\toprule
Audit & $n_{train}$ & $n_{test}$ & $\Delta$RMSE & 95\% CI & Status \\
\midrule
Rolling fold 1 (2004--2006) & 473 & 103 & -0.010 & [-0.039, 0.030] & FAIL \\
Rolling fold 2 (2007--2009) & 561 & 102 & 0.017 & [0.003, 0.028] & FAIL \\
Rolling fold 3 (2010--2012) & 663 & 102 & -0.015 & [-0.090, 0.017] & FAIL \\
Rolling fold 4 (2013--2015) & 769 & 105 & -0.035 & [-0.053, -0.025] & PASS \\
Prefix 25\% through 1996 & 223 & 293 & 0.030 & [0.020, 0.040] & FAIL \\
Prefix 50\% through 2002 & 425 & 313 & -0.011 & [-0.028, 0.007] & FAIL \\
Prefix 75\% through 2009 & 663 & 313 & -0.016 & [-0.036, 0.001] & FAIL \\
Prefix 100\% through 2015 & 879 & 333 & -0.005 & [-0.019, 0.009] & FAIL \\
\bottomrule
\end{tabular}
}
\end{table*}

\subsection{Audit Traceability and Interpretation Limits}
Table~\ref{tab:trace} maps each numerical claim family to a concrete artifact.
This is intentionally more than a code-release statement. E1 records the temporal
split and train-only target construction; E2 records the entire validation grid,
fixed seed outcomes, and tie rule; E3 stores every sampled calendar-year block;
and E4 records every tail prediction, metric, and paired interval. The remaining
checks preserve the temporal, feature, crop, spatial, target-construction, and
baseline evidence required to reproduce the stop decision. Consequently, a reader
can distinguish an empirical model result from a narrative inference about weather.

\begin{table}[t]
\centering
\caption{E1--E10 numerical traceability. VERIFIED means the evidence artifact
was regenerated and checked; it does not mean a scientific gate passed.}
\label{tab:trace}
\scriptsize
\resizebox{\columnwidth}{!}{\begin{tabular}{lll}
\toprule
ID & Artifact & Status \\
\midrule
E1 & split\_manifest.csv & CHECKED \\
E2 & selected\_config.json & CHECKED \\
E3 & figure2\_three\_comparisons.json & CHECKED \\
E4 & tail\_metrics\_by\_threshold.csv & CHECKED \\
E5 & temporal\_and\_capacity\_audits.csv & CHECKED \\
E6 & stage\_feature\_sensitivity.csv & CHECKED \\
E7 & group\_macro\_metrics.csv & CHECKED \\
E8 & retrospective\_target\_comparison.csv & CHECKED \\
E9 & baseline\_vector\_hashes.csv & CHECKED \\
E10 & final\_pdf\_numerical\_crosscheck.json & CHECKED \\
\bottomrule
\end{tabular}
}
\end{table}

Several limits should shape interpretation. State-representative weather
coordinates can miss within-state exposure; the weather covariates are correlated;
and this panel lacks soil, irrigation, cultivar, management, and phenological
measurements. The response variable is a detrended state aggregate rather than a
farm-level loss. The observed gate failure may therefore reflect a mismatch among
target, resolution, features, and model family, not an absence of physical weather
effects. Conversely, these limitations do not justify relaxing the gate after
observing the locked test. They identify inputs for a future, separately specified
study: richer spatial covariates, prospectively defined phenology windows, and an
independent evaluation period.

\subsection{Retrospective Target Sensitivity}
Full-series detrending uses later yields to construct earlier residuals. It is
therefore a retrospective diagnostic, not a prospective target. Relative to the
train-only construction, the $z<-1$ final-period set has Jaccard overlap 0.605
with year-block 95\% interval [0.500, 0.675]; 21 prospective events are not
retrospective events and 13 retrospective events are not prospective events
(Fig.~\ref{fig:leakage}). A coherent retrospective explanation would still not
repair the prospective fidelity failure.

\begin{figure}[t]
  \centering
  \includegraphics[width=\linewidth]{figure_retrospective_v2.png}
  \caption{\textbf{RETROSPECTIVE ONLY.} Full-series detrending uses future yields.
  Bars show exact event counts. Transition labels mean Neither, Train-only only,
  Full-series only, and Both; this is not prospective forecast evidence.}
  \label{fig:leakage}
\end{figure}

\section{Discussion}
\subsection{Claim Scope and Multiplicity}
The audit distinguishes primary decisions from descriptive measurements before
examining the locked period. Gate A is a conjunction on the selected
weather-only model, overall paired error, and the $z<-1$ tail; primary Gate B1 is
the independent incremental-weather-value comparison that cannot alter selection.
Gate B2 is sensitivity only and also cannot alter selection. The two
more severe tails, alternative minimum-history rules, Huber detrending,
standardized targets, rolling folds, prefix learning curves, feature-stage
proxies, crop/state checks, and additional model families are sensitivity or
diagnostic analyses. They provide uncertainty about the protocol boundary, but
they are not multiple opportunities to declare the primary claim successful.
No sensitivity replaces the validation-locked configuration, primary tail, or
test population. This scope statement is important because an isolated favorable
interval, such as a later-history or one rolling-origin result, is expected among
many correlated analyses and does not repair the pre-specified gate failure.

\subsection{Reproducibility and Traceability}
The audit records the source-input hashes, target construction, locked
configuration, fixed seeds, row identifiers, predictions, paired bootstrap draws,
and table inputs. Each numerical claim is traceable to an evidence file; VERIFIED
means regenerated and checked, never scientifically favorable. An anonymized
artifact is prepared for release where permitted by the venue and source licenses.

Traceability is especially important here because several quantities that can look
interchangeable are not. The raw residual in t ha$^{-1}$ is the prediction target;
the training-only standardized residual defines the event threshold; and a cluster
count is the unit sampled for uncertainty, not the number of test rows. The released
records retain the row-aligned target and prediction vectors, zero-residual baseline,
calendar-year bootstrap draws, and configuration hash. A reader can therefore
recompute a reported difference in RMSE, inspect its confidence interval, and check
that both sides of a comparison use identical locked rows. This design also makes a
failed check informative: an artifact can be complete and a gate can still fail.

The selection boundary is similarly explicit. Candidate models, feature families,
and seeds are evaluated only in validation, after which one weather-only ExtraTrees
configuration is locked. The final period is then used for the paired Gate A/B audit
rather than for a second search. Alternative target constructions, histories,
resampling units, and model families are preserved with their roles in the evidence
table. A favorable sensitivity is consequently visible but cannot be silently
promoted into the primary conclusion. This is a practical guard against multiple
correlated analyses creating an attractive but unsupported narrative.

The event metrics follow the same boundary. Precision, recall, F1, balanced
accuracy, PR-AUC, and ROC-AUC are calculated over every locked row using the fixed
zero score threshold, with the confusion-matrix counts retained in the artifact.
They describe discrimination under a fixed operational rule, whereas Gate A asks a
stricter question: whether the selected model improves paired error and recovers
the primary tail beyond the specified null-aware checks. No collection of
diagnostics substitutes for that conjunction. This distinction is why the paper
reports a descriptive event table while withholding event-driver interpretation.

\subsection{Limitations and External Validity}
State-representative weather can miss within-state exposure, weather covariates are
correlated, and the panel lacks soil, irrigation, cultivar, management, and
phenology data. This negative result limits the present target, resolution,
features, and model family; it does not establish that weather has no yield effect.
A future study needs prospectively versioned inputs, finer exposure measurements,
and an independent future test interval.

The temporal and population sensitivities also limit external validity. Three of
four rolling folds have negative point estimates, but only one has a confidence
interval entirely below zero; the direction and uncertainty are therefore not stable
enough for a general predictive claim. Raising the minimum history from three to
eight or ten years changes eligibility from 333 to 313 locked rows. Its favorable
interval is useful evidence that series history matters, but it estimates a different
population and is labeled as such. Full-series detrending changes target membership
by using future yields and remains retrospective only. These boundaries identify
what should be pre-specified and externally validated before a later study revisits
the predictive or explanatory question.

\section{Synthetic and External-Domain Evidence}
The stop-gate policy is also evaluated where ground truth is available. A
deterministic synthetic temporal benchmark covers no, weak, moderate, and strong
signal; correlated features; omitted confounding; temporal and geographic shift;
leakage; small samples; imbalanced tails; measurement error; and spatial-resolution
mismatch. Ungated explanations permit attribution in every scenario, including
invalid regimes. The full gate reduces simulated false permission from 28.6\% to
0\% while permitting explanation in six valid-signal scenarios. Thus the gate is
not an always-abstain rule; it is a conditional permission rule.

An external, non-agricultural stress test uses public daily PJM electricity-demand
records. Calendar-only features are fit on January--September 2024 and evaluated
on a locked October--December period. Adding the day-ahead demand forecast improves
locked RMSE from 227,417 to 77,521 MWh, with paired bootstrap $\Delta$RMSE
95\% CI $[-186,037,-119,309]$. The corresponding feature-group gate passes, so
permutation importance is made available for this external forecasting task. This
result establishes that the procedure can permit interpretation when its
pre-specified predictive and incremental-value conditions are met. It is neither a
causal claim nor evidence of agricultural transfer.

\section{Conclusion}
This paper evaluates whether an interpretable weather-only model can support
substantive claims about below-trend U.S. crop-yield events. The model and feature
family were selected before a locked 2016--2025 test. Gate A requires paired
predictive fidelity and null-aware recovery of the primary $z<-1$ tail; primary
Gate B1 requires incremental weather value of the selected weather-only model over
metadata. Gate A and Gate B1 fail. Accordingly,
the permitted conclusion is procedural and negative: this fitted model may be
described, but its explanations are not evidence of observed event drivers or
weather-specific effects in this panel. The reusable lesson is to make data
availability and target construction explicit, lock selection before testing,
compare same-row baselines with paired uncertainty, distinguish predictive
fidelity from feature value, and stop interpretation when a pre-specified gate
fails. The audit artifacts preserve the numerical path to that decision.

\bibliographystyle{IEEEtran}
\bibliography{references}
\end{document}
